% Nouveau EGA V (cote 25): archive cote 25, site page 133 (PDF page 134).
% typescript. Draft compared: none.
% AI-transcribed by Claude Opus 5.5 (claude-opus-5-5, adaptive thinking, effort high), Message Batches API, batch ega25b (msgbatch_01A7n5YuXubUPNK1zsepi8RG), 2026-11-13; not yet checked by eye.
\scanpage{25/133}
\mspage{\hand{36}}

un sous-k-espace vectorielEde l'espace des  fonctions  méromorphes  
sur X, tel que pour tout $\hand{\phi} \hand{\in} E-(0)$, $\hand{\phi}$ soit régulier 
  , i.e. div($\hand{\phi}$) est défini 
\hand{\unsure{,|)}} et que $\hand{\triangle}$ soit l'ensemble des div($\hand{\phi}$) pour $\hand{\phi} \hand{\in} E-(0)$.
c) Soient E, E' deux sous-k-espaces vectoriels de l'espace des
fonctions rationelles strictes  sur X,  satisfaisant la condition
envisagée dans b), alors les ensembles de diviseurs $\hand{\triangle}$, $\hand{\triangle}'$ qu'ils
définissent respectivement sont égaux sss il existe une pseudo- fonction
e  régulière $\hand{\phi}$ sur X, telle que $E' = \hand{\phi} E$ . Si $E\hand{\neq}(0)$(i.e. $\triangle \neq \emptyset$), une
telle $\hand{\phi}$ est déterminée modulo multiplication par un élément c de $k^{\hand{*}}$.

La démonstration est un exercice facile, enutilisant \hand{4}.1. , et je
me dispense d'en écrire la démonstration, sauf réclamation de ta part. Il
\placedmargin{||} me sembled'ailleurs que \hand{4}.4. viendrait avantageusement dès avant \hand{4}.3.,
étanttechniquement nettement plus trivial. Noter d'ailleurs que lorsque
X est géom intègre, la condition sur E énoncée dans b) devient \uline{vide}.
La restriction mise au début de b) tient au fait qu'il se pourrait autre-%
ment que la condition énoncée pour b) ne soit plus vérifiée après passage
à la clôture algébrique de k (il est facile d'en donner des exemples,
en toute caractéristique, et même lorsque k est séparablement clos en
c p \hand{$>$} 0). Pour bien faire, on devrait énoncer b) sans condition supplé-%
mentaire sur k ou X, mais en énoncant la condition sur E en passant à la
clôture algébrique de k (et signaler que si X est géom intègre, \unsure{oette|cette}
condition devient vide). \hand{-- Par abus de langage, un}
\hand{\uline{ens.} $\triangle$ de diviseurs de la forme \unsure{Ens|En}(D) sera}
\hand{encore parfois appelé système linéaire de diviseurs}
\hand{sur X ;}

\mspage{\hand{133}}
